% !TEX root = ../../../main.tex
% Appendix entry: catalogue.cluster -- "Hierarchical cluster" (Appendix A).
\entryfield{Algorithm}
\code{_run} in \code{Adapters/catalogue_cluster.py} (\code{_cluster_window_set}) calls
\code{preprocess_window_matrix} and \code{cluster_window_matrix} in
\code{Working/Catalogue/dendrogram/dendrogram_cluster.py}:
\begin{enumerate}
\item Drop the bookkeeping columns (\code{BOOKKEEPING_COLUMNS} in
  \code{Working/types/windowset.py}: \code{split}, \code{event_index}, \code{anchor_offset},
  \code{overlaps_next}, \code{span_score}, family and cycle columns, and others). A window whose
  remaining features are all NaN is left out ($-1$) and counted; $k$ larger than the windows left
  is refused.
\item \code{preprocess_window_matrix}, with its defaults: set aside the metadata columns
  \code{category}, \code{fusion_pred_v1}, \code{fusion_pred_v1_error}; drop non-numeric columns,
  all-NaN columns, columns more than 50\% NaN, constant columns, and near-constant columns (standard
  deviation below $10^{-8}$ of the column's magnitude); drop the later column of any pair with
  $|r| > 0.99$ (skipped above 2000 columns); impute remaining NaNs with the column median; z-score
  every column (\code{StandardScaler}) over \emph{all} windows.
\item \code{cluster_window_matrix(method=linkage, n_clusters=k)}: for Ward, \code{scipy}'s
  \code{linkage(X, "ward")} on the scaled matrix \citep{ward1963hierarchical}; for the others, on
  condensed Euclidean distances. Optimal leaf ordering below 2000 windows. Flat clusters by
  \code{fcluster(Z, k, "maxclust")}, labelled 1..$k$. Silhouette on the scaled features
  \citep{rousseeuw1987silhouettes}, sampled above a size cap
  (\code{_SILHOUETTE_SAMPLE_MAX_N}); cophenetic correlation when the distance array is affordable
  (under 400\,MB), else \code{None}.
\end{enumerate}
$k$ is never chosen by the block: the code says the selection criterion is an open research
decision, so it stays a setting.

\entryfield{Assumptions}
\begin{itemize}
\item \emph{Units.} The z-scoring makes each column unitless; a column's unit only matters through
  the near-constant filter.
\item \emph{Missing values.} All-NaN windows are left out; other NaNs are median-imputed, which
  pulls those windows towards the centre of the cloud.
\item \emph{Features, not bookkeeping.} Only the columns on \code{BOOKKEEPING_COLUMNS} (and the three
  metadata names) are excluded. Any other numeric column is clustered on.
\item \emph{No roles.} The \code{split} column is dropped as bookkeeping, not used: every window is
  scaled and clustered together.
\end{itemize}

\entryfield{Outputs}
A \texttt{Grouping}: 1..$k$ per clustered window, $-1$ for a left-out one. \code{meta}:
\code{linkage}, \code{k}, \code{cluster_counts}, \code{n_windows}, \code{n_clustered},
\code{n_left_out} and its reason, \code{summary}, \code{n_features}, \code{silhouette},
\code{cophenetic_r}, \code{feature_names} (the columns actually used), \code{correlation_filter}
(applied or skipped). The block has a pre-run readout of cluster sizes (\code{_derive}). It
persists: the labels are written as a CSV (window index, label) under
\code{DATA/derived/groupings/} and registered as an \code{artifacts(kind='csv')} row.

\entryfield{Cost}
Ward on the feature matrix needs no explicit distance array, but scipy's linkage is still
$O(n^2)$ time; other linkages, optimal ordering and the cophenetic correlation need $n(n-1)/2$
doubles and are refused above 2048\,MB (about 23,000 windows). The estimate is always
\code{None} (unknown): the cost depends on the window count, which the estimator cannot see.

\entryfield{Failure modes}
\begin{itemize}
\item \emph{$k$ clusters from anything.} The cut always yields $k$ groups; read the silhouette.
\item \emph{Test windows in the clustering.} When the windows carry a split, test windows shape
  the scaling and the tree. A classifier trained on these labels has seen the test windows'
  cluster structure. The pool route (\blockref{catalogue.shape_cluster}) avoids this.
\item \emph{A pool fed directly.} The \code{pool_*} metadata columns of
  \blockref{preprocessing.window_pool} (recording id, channel, start, length, rate, scale, set id,
  stretch bounds) are not on the bookkeeping list, so they would be clustered as features; only
  the text columns are dropped. Other readers refuse a pool (MiniRocket, multiscale entropy); this
  one does not.
\item \emph{Chaining} with \texttt{single} or \texttt{average} linkage gives one big cluster and
  singletons; the code logs it (\code{_check_cluster_imbalance}).
\item \emph{CNN and RF columns} from \blockref{preprocessing.window_matrix} carry the human labels
  into the clustering.
\end{itemize}

\entryfield{Validation}
\code{tests/test_grouping_adapter.py} recovers three obvious synthetic clusters by membership.
Nothing tests a choice of $k$ or linkage on real windows.

\entryfield{References}
\citet{ward1963hierarchical}; \citet{rousseeuw1987silhouettes}. Research question 1:
\code{docs/rq_roundA/RQ1-cluster-vs-manual-labels.md}.

\entryfield{Tests}
\code{tests/test_grouping_adapter.py}: signature; linkage and $k$ are settings; the readout; three
clusters recovered; round trip to disk; meta carries silhouette, cophenetic $r$ and feature names;
an empty set refused. \code{tests/test_all_nan_feature_rows.py}: all-NaN rows left out, $-1$,
counted. \code{tests/test_correlation_filter_ceiling.py}: the 2000-column ceiling.
\code{tests/test_adapter_sliding_windows.py} and \code{tests/test_adapter_event_windows.py}: the
split and the event bookkeeping are not used as features.
